% -*- coding: utf-8 -*-
% 欢迎使用 YangTex —— 中英文实时编译示例
% 编译引擎：xelatex（中文请始终使用 xelatex / lualatex）
\documentclass[11pt,a4paper,fontset=fandol]{ctexart}

\usepackage{amsmath,amssymb,amsthm,mathtools}
\usepackage{geometry}
\usepackage{hyperref}
\geometry{margin=2.6cm}
\hypersetup{colorlinks=true,linkcolor=blue,citecolor=blue,urlcolor=blue}

\newtheorem{theorem}{定理}[section]
\newtheorem{lemma}[theorem]{引理}
\newtheorem{definition}[theorem]{定义}
\newcommand{\R}{\mathbb{R}}
\newcommand{\Om}{\Omega}
\DeclareMathOperator{\divg}{div}

\title{\textbf{欢迎使用 YangTex}\\[4pt]
       \large 轻量 \quad 实时 \quad 中英文 \quad 精确双向同步}
\author{杨建伟 \quad 数学与统计学院}
\date{\today}

\begin{document}
\maketitle

\begin{abstract}
\noindent YangTex 是一个本地运行的轻量级 \LaTeX{} 写作环境：输入即编译，中英文混排开箱即用，
内置 AMS 等数学期刊模板，并通过 SyncTeX 实现源码与 PDF 之间的\textbf{逐行双向定位}。
本页即是实时编译的成品，试着在左侧修改任意一行，右侧会立即更新。
\end{abstract}

\section{开始写作}

把光标放在这一行，右侧预览会自动滚动到对应位置；单击右侧 PDF 的任意文字，
左侧编辑器就会跳回精确的源码行。这就是双向同步（SyncTeX）。

中文段落与 English paragraph 可以自然混排，公式同样即时排版：
\begin{equation}\label{eq:model}
  -\Delta u + |u|^{p-2}u = f \quad \text{in } \Om\subset\R^n,
  \qquad u|_{\partial\Om}=0 .
\end{equation}

\begin{theorem}[能量估计]
设 $f\in L^2(\Om)$ 且 $u\in H^1_0(\Om)$ 是方程 \eqref{eq:model} 的弱解，
则存在仅依赖于 $\Om$ 与 $p$ 的常数 $C>0$，使得
\begin{equation}
  \|\nabla u\|_{L^2(\Om)} \le C\,\|f\|_{L^2(\Om)} .
\end{equation}
\end{theorem}

\begin{proof}
在 \eqref{eq:model} 两端乘以 $u$ 并在 $\Om$ 上积分，由分部积分与
Poincaré 不等式即得结论。
\end{proof}

\section{常用功能  追着}
\begin{itemize}
  \item \textbf{自动编译}：停止输入片刻后自动编译，或按 \texttt{Cmd/Ctrl + Enter} 手动编译。
  \item \textbf{模板库}：顶部「模板」可一键套用 amsart、amsproc、amsbook、Elsevier、IEEE、
        Springer、SIAM、中文 ctex、Beamer 等；把自己的 \texttt{.tex} 放进 \texttt{templates/}
        目录即可自动加入模板列表。
  \item \textbf{AI 辅助}：右侧 AI 面板可解释公式、润色文字、生成代码、诊断编译错误，
        接入免费大模型，无需任何 API Key 设置。
  \item \textbf{工具面板}：顶部「\texttt{工具面板}」按宏命令、运算符、定界符、环境、
        函数、字体、表格等分类一键插入代码片段；选中文字时会自动包裹进命令参数。
  \item \textbf{新建项目}：每个项目独立目录，PDF 与源文件保存在同一处。
\end{itemize}

\section{A Small English Section}
This paragraph is written in English to demonstrate that YangTex handles
both Chinese and English in a single document. Mathematical environments
from \textsc{AMS} packages are available:
\begin{align}
  \frac{\partial u}{\partial t} - \divg(A(x)\nabla u) &= 0, \\
  \|u\|_{L^\infty(0,T;L^2(\Om))} + \|\nabla u\|_{L^2(Q_T)} &\le C.
\end{align}

\bigskip
\noindent\rule{\linewidth}{0.4pt}
\begin{center}
\small 祝写作顺利。\quad \textit{Happy \TeX ing with YangTex!}
\end{center}

\end{document}
